\chapter{引言与数学基础}
\intro{本章介绍数学物理方程的基本概念和必要的数学工具。}

\begin{mydefinition}{三大典型方程}
\begin{enumerate}
	\item \textbf{波动方程}（双曲型）：$\displaystyle\frac{\partial^2 u}{\partial t^2}=a^2\lapl u$，描述振动、声波、电磁波。
	\item \textbf{热传导方程}（抛物型）：$\displaystyle\frac{\partial u}{\partial t}=k\lapl u$，描述热扩散、物质扩散。
	\item \textbf{拉普拉斯方程}（椭圆型）：$\lapl u=0$，描述稳态温度、静电场、引力场。
\end{enumerate}
\end{mydefinition}

\section{求解方法概述}
\begin{enumerate}
	\item 分离变量法：将 PDE 转化为 ODE 求解。
	\item 积分变换法：利用傅里叶变换、拉普拉斯变换。
	\item 格林函数法：利用格林函数构造特解。
	\item 变分法：转化为泛函极值问题。
	\item 数值解法：有限差分法、有限元法。
\end{enumerate}

\section{向量分析基础}

\begin{mydefinition}{梯度、散度与旋度}
设 $f$ 为标量场，$\vect{A}=(A_x,A_y,A_z)$ 为矢量场：
\begin{align*}
	\grad f&=\frac{\partial f}{\partial x}\vect{e}_x+\frac{\partial f}{\partial y}\vect{e}_y+\frac{\partial f}{\partial z}\vect{e}_z\\
	\divg\vect{A}&=\frac{\partial A_x}{\partial x}+\frac{\partial A_y}{\partial y}+\frac{\partial A_z}{\partial z}\\
	\curl\vect{A}&=\left(\frac{\partial A_z}{\partial y}-\frac{\partial A_y}{\partial z}\right)\vect{e}_x+\left(\frac{\partial A_x}{\partial z}-\frac{\partial A_z}{\partial x}\right)\vect{e}_y+\left(\frac{\partial A_y}{\partial x}-\frac{\partial A_x}{\partial y}\right)\vect{e}_z
\end{align*}
\end{mydefinition}

\begin{theorem}[格林第一恒等式]
\[
\iiint_V(\phi\lapl\psi+\grad\phi\cdot\grad\psi)\,dV=\oiint_S\phi\frac{\partial\psi}{\partial n}\,dS
\]
\end{theorem}

\begin{theorem}[格林第二恒等式]
\[
\iiint_V(\phi\lapl\psi-\psi\lapl\phi)\,dV=\oiint_S\left(\phi\frac{\partial\psi}{\partial n}-\psi\frac{\partial\phi}{\partial n}\right)dS
\]
\end{theorem}

\begin{remark}
格林恒等式是推导互易定理、唯一性定理和格林函数法的基础工具。
\end{remark}

\section{傅里叶级数}

\begin{mydefinition}{傅里叶级数}
$f(x)$ 在 $[-\pi,\pi]$ 上的展开：
\[
f(x)=\frac{a_0}{2}+\sum_{n=1}^{\infty}(a_n\cos nx+b_n\sin nx)
\]
其中
\[
a_n=\frac{1}{\pi}\int_{-\pi}^{\pi}f(x)\cos nx\,dx,\quad b_n=\frac{1}{\pi}\int_{-\pi}^{\pi}f(x)\sin nx\,dx
\]
\end{mydefinition}

\begin{figure}[htbp]
\centering
\begin{tikzpicture}
\draw[->,thick] (-2,0) -- (2,0) node[right]{$x$};
\draw[->,thick] (0,-1.5) -- (0,2.5) node[above]{$y$};
\draw[->,thick] (0,0) -- (1.5,1.5) node[above left]{$\vect{A}$};
\draw[dashed,blue] (1.5,0) -- (1.5,1.5);
\draw[dashed,red] (0,1.5) -- (1.5,1.5);
\node[below right] at (1.5,0) {$A_x$};
\node[left] at (0,1.5) {$A_y$};
\end{tikzpicture}
\caption{向量分解}
\end{figure}
